\documentclass[10pt,twocolumn]{article}
\usepackage[margin=0.85in]{geometry}
\usepackage{times}
\usepackage{hyperref}
\usepackage{xcolor}
\usepackage{listings}
\usepackage{amsmath}

\hypersetup{colorlinks=true, linkcolor=black, citecolor=black, urlcolor={blue!60!black}}

\lstset{basicstyle=\ttfamily\footnotesize, breaklines=true, columns=fullflexible,
  frame=single, framesep=4pt, xleftmargin=4pt, aboveskip=6pt, belowskip=6pt}

\title{\textbf{PriceRight: An Accountability Layer for\\ x402 Agent Commerce on Base}}
\author{Dipankar Sarkar}
\date{2026}

\begin{document}
\maketitle

\begin{abstract}
Autonomous agents are starting to transact with each other on Base, where x402 makes an
agent-to-agent USDC payment a single HTTP round trip, but an agent that is paid to produce a
judgement has no skin in the game: if it answers wrong, it keeps its fee and moves on.
PriceRight closes that gap by making the right to answer a paid resource and the answer a
collateralised commitment. The claim endpoint replies \texttt{402 Payment Required} with
x402 payment requirements; the resolver signs an EIP-3009 \texttt{TransferWithAuthorization}
and presents it in an \texttt{X-PAYMENT} header, and the arena settles that authorisation
on-chain, so the fee moves by signature rather than by a prior allowance and is bound to one
task by a nonce that commits to the arena address and the chain id. In the same call the
resolver's collateral is custodied per task in a vault keyed by its ERC-8004 identity, which
is a real ERC-721. Settlement is a pure function of the committed verdict and a truth the
poster hash-committed in advance: correct returns the fee, the bounty and the collateral and
writes a score of 100 into the ERC-8004 reputation registry; wrong moves the collateral to
the poster and writes a zero. Commit and reveal deadlines with permissionless timeouts mean
neither party can freeze the escrow, and no terminal path leaves value stranded. We report
46 Solidity tests, 53 Python tests, a cross-language pin in which a signature produced by
this paper's pure-Python secp256k1 is verified by Solidity's \texttt{ecrecover} while the
Python suite regenerates the same vector from the Solidity constants, and a recorded local
devnet run whose transaction hashes are reproduced in the interface.
\end{abstract}

\section{Problem}
Agent-to-agent commerce assumes an agent that takes payment for work will do the work
well, but a fee alone creates no accountability. If a resolver agent is paid to answer a
subjective yes/no question (``did this event happen?'', ``did this project ship?'') and
answers wrong, a naive market simply pays it and forgets. The result is an economy where
confidently wrong agents are indistinguishable, ex ante, from careful ones, and where a
buyer of judgements cannot price the risk of a bad answer.

Two ingredients are missing. First, a payment rail that gates access to the task so that
answering is never free~\cite{x402,rfc9110}. Second, a portable, on-chain identity that
carries reputation \emph{and} bonded value, so a wrong answer costs the agent something it
cannot quietly reset~\cite{erc8004}. Proof-of-stake systems have long used bonded stake
plus slashing to make misbehaviour expensive~\cite{bentov2016snowwhite,nakamoto2008bitcoin};
PriceRight applies the same economic logic to the outputs of autonomous agents.

\section{Related Work}
\textbf{ERC-8004 Trustless Agents}~\cite{erc8004} specifies an Identity registry (an agent
as an ERC-721~\cite{eip721} whose token id is the agent id), a Reputation registry in which
feedback is permissionless and readers filter by client, and a Validation registry for
third-party attestation. PriceRight implements all three and adds the economics the
standard deliberately omits: a collateral vault keyed by agent id, so a score of zero is
accompanied by value that actually moved. Because reputation is keyed by agent id rather
than by address, a transferred identity carries its record with it. The fee token follows
ERC-20~\cite{eip20} and additionally implements EIP-3009 transfer authorisations, as USDC
does.

\textbf{x402}~\cite{x402} is an HTTP-native payment standard built on the 402 Payment
Required status code~\cite{rfc9110}: a resource answers 402 with an \texttt{accepts} array
of payment requirements, the client returns a payment payload in an \texttt{X-PAYMENT}
header, and a facilitator verifies and settles it. On EVM chains the \texttt{exact} scheme
settles by EIP-3009 over USDC on Base, which is what PriceRight implements: the right to claim and be scored
on a task is the paid resource, and the arena is the party that submits the payer's signed
authorisation.

\textbf{Bonded stake and slashing} are the core accountability mechanism of proof-of-stake
consensus~\cite{bentov2016snowwhite}, where a validator that signs conflicting messages
loses stake. Adversarial value-extraction on public chains~\cite{daian2020flash} motivates
a deterministic, commit-based settlement so the outcome cannot be steered after the fact.
Content integrity uses keccak-256~\cite{keccak2011}, the hash Ethereum exposes as
\texttt{keccak256}~\cite{buterin2014ethereum}.

\section{Design}
PriceRight is six contracts and a resolver agent.

\paragraph{ERC-8004 registries.} \texttt{IdentityRegistry} is an ERC-721 whose token id is
the agent id; \texttt{supportsInterface} reports \texttt{0x80ac58cd} and
\texttt{0x5b5e139f}, and identities transfer and approve like any other token.
\texttt{ReputationRegistry} exposes \texttt{giveFeedback} to anyone and
\texttt{getSummary(agentId, clients, tag)} to readers, who decide whose feedback counts;
the arena writes as an ordinary client, not as an administrator.
\texttt{ValidationRegistry} records a request and, later, a response.

\paragraph{AgentStakeVault.} ERC-8004 says nothing about money, so the economics live in an
extension. Collateral is real ERC-20 held by the vault under a key the arena derives from
$(\text{arena}, \text{taskId})$ and tagged with the agent id it backs. Every bond ends in
exactly one of two terminal states, and both move tokens: \texttt{release} to a named
address, or \texttt{slash} to a named beneficiary. There is no path that leaves collateral
in the vault, which the suite asserts by draining the system after each settlement.

\paragraph{AgentArena.} A poster calls \texttt{postTask} with a bounty and a salted
commitment $c = \mathrm{keccak256}(\mathrm{uint8}(t) \,\|\, s)$, fixing the truth $t$ before
any resolver can commit. A resolver calls \texttt{claimTask} with an x402 \texttt{exact}
payment: the EIP-3009 fields plus $(v,r,s)$. The arena requires that the payer is the
caller, that the value equals the posted fee, and that the nonce equals
\[
\mathrm{keccak256}\big(\textsc{scope} \,\|\, \text{arena} \,\|\, \text{chainId} \,\|\,
\text{taskId} \,\|\, \text{agentId}\big),
\]
then calls \texttt{transferWithAuthorization} and checks the balance actually moved. That
nonce is the binding: an authorisation signed for one task cannot be replayed against
another task, another arena or another chain, and the token's own nonce map makes it
single-use. Collateral is bonded in the same call. The resolver then calls
\texttt{commitVerdict}, which stores $\mathrm{keccak256}$ of its reasoning and files a
validation request whose request hash is derived from it, before the truth is revealed, so
the attestation cannot be back-dated. Finally the poster calls \texttt{settle}, revealing
$(t,s)$; the contract checks the reveal reproduces $c$ and compares:
\[
\text{correct} \Rightarrow \text{fee} + \text{bounty} + \text{collateral to the resolver},\
\text{feedback } 100
\]
\[
\text{wrong} \Rightarrow \text{fee} + \text{bounty} + \text{collateral to the poster},\
\text{feedback } 0.
\]

\paragraph{Liveness.} A claim sets a commit deadline and a commit sets a reveal deadline.
After either passes, \emph{anyone} may call the corresponding timeout. A resolver that
never commits forfeits its fee to the poster but keeps its collateral, since it gave no
wrong verdict; a poster that never reveals pays the resolver the fee, the bounty and its
collateral back, which makes withholding a reveal strictly worse than revealing whatever
the truth turns out to be. Payouts are credits withdrawn by pull, so one unlucky transfer
cannot brick a settlement.

\paragraph{Two backends, one interface.} The agent talks to a chain object.
\texttt{JsonRpcChain} builds, RLP-encodes, signs and broadcasts EIP-1559 transactions and
reads state with \texttt{eth\_call}. \texttt{InMemoryChain} executes the same state machine
in process, including EIP-3009 verification against a real EIP-712 digest with real
signature recovery, so a forged or replayed authorisation fails there exactly as it fails
in the token. Configuration selects between them in one place, and a half-configured
environment raises rather than quietly running the mirror.

\paragraph{The resolver does not know the answer.} The poster holds $(t,s)$; the resolver
is handed a claim and an evidence bundle. \texttt{EvidencePolicy} parses a threshold claim
such as \texttt{ETH/USD >= 4000 @ block 21451200}, keeps the readings that match both the
symbol and the block, takes their median and compares. When the evidence does not cover the
claim it returns an abstention and the agent declines to pay at all. The contrasting agent,
\texttt{StaleCachePolicy}, evaluates the same comparison against a week-old cached quote.
It is not adversarial and it is not rigged to lose: given a threshold its cache clears, it
wins.

\section{The Hero}
Two agents, one arena, one rule, on a live node. Both pay the same x402 fee. The first
reads three oracle readings, commits \emph{Yes}, and at settlement is credited the fee, the
bounty and its collateral, with a score of 100 written to the reputation registry. The
second answers from its cache, commits \emph{No}, and at settlement its collateral is moved
to the poster and a zero is written against its identity. The difference between them is
diligence, not honesty, which is the failure mode a market of paid agents actually has.

What makes it a hero rather than an animation is that a reviewer can leave the room and
check it: \texttt{cast logs} shows the \texttt{Slashed} event, \texttt{cast call} on the
vault returns the slashed amount for that agent id, and the interface displays the
transaction hashes from the same recorded run.

\section{Evaluation}
The contracts compile under solc 0.8.24 with no vendored dependencies and pass 46 Foundry
tests: the lifecycle, both settlement paths, the value-conservation invariant after each,
both permissionless timeouts, the deadline guards, the x402 binding and replay properties,
and ERC-721 and reputation conformance. The Python side passes 53 tests covering the same
state machine, the policy's dependence on evidence, the abstention path, the streamed run
events and every rejection reason the facilitator can return.

Two results are worth separating from the counts. First, the fee genuinely moves by
signature: the resolver holds no allowance for the arena, so
\texttt{test\_fee\_moves\_by\_signature\_not\_by\_allowance} would fail if the payment were
decorative. Second, the two languages are pinned to each other. The Python client signs
with the pure-Python secp256k1 in this repository; a vector it produced is fed to Solidity,
which re-derives the EIP-712 digest and recovers the payer with \texttt{ecrecover}, while
the Python suite reads those constants back out of the Solidity file and regenerates them.
Editing either side to fit the other fails one of the two tests.

\begin{lstlisting}
$ forge test
46 passed; 0 failed
$ pytest tests -q
53 passed
$ ./scripts/devnet.sh
agent #1  evidence     -> CORRECT  +110 tUSD, collateral returned, 100%
agent #2  stale-cache  -> WRONG    50 tUSD collateral seized,        0%
$ cast call $VAULT 'slashedOf(uint256)(uint256)' 2
50
\end{lstlisting}

\section{Limitations}
The system has not yet been deployed to Base Sepolia: no funded testnet key was available in
this environment, so there is no explorer link, and what has actually been executed is the
local devnet run described above; the RPC adapter is network-agnostic and the same code path
runs against Base Sepolia once a funded key is supplied, which is the milestone-1 target. No
hosted CDP facilitator has been called either, because that needs credentials:
the remote rail is exercised against a conformant facilitator served over HTTP by the test
suite, which covers the approved, refused and unreachable cases. Verification can be remote
but settlement cannot, because \texttt{claimTask} consumes the authorisation itself; a
facilitator that broadcast it first would burn the nonce and make the claim impossible.

Ground truth is still supplied by the poster under a commitment. The commitment binds the
poster to one value in advance and the reveal timeout removes the incentive to withhold,
but neither proves the value is real; sourcing truth from an oracle or a dispute game is the
obvious next step. The collateral is a fixed per-task amount rather than a function of
confidence or repeat offences. The resolver is a deterministic evidence parser rather than a
model, which is a choice made for reproducibility in front of a reviewer: a model fits
behind the same \texttt{decide(claim, evidence)} interface, and the arena does not care how a
verdict was formed. Finally, the ERC-8004 registries here are a self-contained
implementation of the standard's structure rather than the reference deployment, and they
omit optional parts of it; \texttt{docs/HONESTY.md} enumerates exactly which.

\bibliographystyle{plain}
\bibliography{references}

\end{document}
